% Отключаем нумерацию страниц для Задания (оно вшивается, но номера на нем не пишутся)
\thispagestyle{empty}

\begin{center}
    Министерство образования Республики Беларусь \\
    Учреждение образования \\
    «БЕЛОРУССКИЙ ГОСУДАРСТВЕННЫЙ УНИВЕРСИТЕТ \\ 
    ИНФОРМАТИКИ И РАДИОЭЛЕКТРОНИКИ» \\
    \vspace{0.5cm}
    Факультет \underline{ИНФОРМАЦИОННЫХ ТЕХНОЛОГИЙ И УПРАВЛЕНИЯ} \\
    Кафедра \underline{СИСТЕМ УПРАВЛЕНИЯ} \\
    
    \vspace{1cm}
    
    \begin{flushright}
        \begin{minipage}{0.5\textwidth}
            УТВЕРЖДАЮ \\
            Заведующий кафедрой \\
            \underline{\hspace{3cm}} И. И. Иванов \\
            «\underline{\hspace{1cm}}» \underline{\hspace{3cm}} 2026 г.
        \end{minipage}
    \end{flushright}

    \vspace{0.8cm}
    {\bfseries ЗАДАНИЕ} \\
    {\bfseries по дипломному проекту (работе)}
\end{center}

\noindent Студент: \underline{Петров Петр Петрович} \hfill Группа: \underline{123456} \\
1. Тема проекта (работы): \underline{Разработка автоматизированной системы мониторинга} \\
\underline{параметров окружающей среды} \\
утверждена приказом по университету от «\underline{\hspace{0.5cm}}» \underline{\hspace{2cm}} 2026 г. № \underline{\hspace{1cm}} \\
2. Срок сдачи студентом законченного проекта (работы): «\underline{\hspace{0.5cm}}» \underline{\hspace{2cm}} 2026 г. \\
3. Исходные данные к проекту (работе): \underline{Технические характеристики датчиков серии DHT,} \\
\underline{микроконтроллер ESP32, протоколы передачи данных MQTT и HTTP.} \\
4. Содержание пояснительной записки (перечень подлежащих разработке вопросов): \\
\underline{Введение, обзор литературы, выбор и обоснование аппаратных средств, разработка} \\
\underline{алгоритмов ПО, экономическое обоснование, раздел охраны труда, заключение.} \\
5. Перечень графического материала (с точным указанием обязательных чертежей и плакатов): \\
\underline{Схема электрическая принципиальная (А1), Алгоритм работы системы (А1),} \\
\underline{Схема данных (А1), Технико-экономические показатели (А1).}

\vspace{0.5cm}

\begin{center}
    {\bfseries КАЛЕНДАРНЫЙ ПЛАН}
\end{center}
\noindent (с указанием сроков выполнения отдельных этапов)
\begin{center}
\begin{tabular}{|p{10cm}|p{3cm}|p{2cm}|}
\hline
Наименование этапов дипломного проекта (работы) & Срок выполнения & Примечание \\ \hline
Изучение литературы и ТЗ & 15.03.2026 & 10\% \\ \hline
Разработка аппаратной части & 01.04.2026 & 30\% \\ \hline
Разработка программного обеспечения & 01.05.2026 & 60\% \\ \hline
Экономика и охрана труда & 20.05.2026 & 85\% \\ \hline
Оформление ПЗ и графического материала & 05.06.2026 & 100\% \\ \hline
\end{tabular}
\end{center}

\vspace{0.5cm}
\noindent Консультанты по проекту (работе): \\
\noindent По разделу экономики: \underline{\hspace{4cm}} /Е. Е. Экономистов/ \\
\noindent По разделу охраны труда: \underline{\hspace{4cm}} /Б. Б. Безопасный/ \\
\noindent Дата выдачи задания: «\underline{\hspace{0.5cm}}» \underline{\hspace{2cm}} 2026 г.

\vspace{0.5cm}
\noindent Руководитель \underline{\hspace{4cm}} /С. С. Сидоров/ \\
\noindent Студент \underline{\hspace{4cm}} /П. П. Петров/

\newpage